\section{Conclusion}
\label{sec:conclusion}

This work frames response personalization as a portable, contextual, transparent,
and user-governed policy-selection problem rather than undifferentiated long-term
memory. It contributes a typed representation, behavioral-conformance definition,
benchmark, executable five-condition harness, local-first product, and blinded
human-evaluation path.

Across 128 real generations from two local model families, the same profile and
selection pipeline executed without failures. A compact compiler saved 19 input
tokens per prompt relative to a short history baseline, but produced about 99 more
output tokens and added 3.12 seconds of latency; its exploratory reference coverage
was not distinguishable in this small sample. Thus the measured contribution is
an auditable systems artifact and a falsifiable efficiency trade-off, not evidence
that personalization improves subjective quality. Establishing that claim requires
the preregistered target-user study and independent correctness evaluation.
